\begin{explanationbox}
	\textbf{\textcolor{CExplanation}{一句话直觉}}
	定比分点公式刻画了线段上（或延长线上）一点按给定比例分配两端点坐标的加权平均关系。

	\medskip
	\textbf{\textcolor{CExplanation}{核心拆解}}
	\begin{description}[style=nextline, leftmargin=2.8em, labelsep=0.5em, itemsep=0.45em]
		\item[\textbf{要点一（向量等式）：}]
			点 $P$ 满足 $\overrightarrow{P_1P} = \lambda \overrightarrow{PP_2}$，该式同时确定了 $P$ 在 $P_1P_2$ 直线上的位置和方向。
		\item[\textbf{要点二（加权平均）：}]
			坐标公式 $x=\frac{x_1+\lambda x_2}{1+\lambda}$ 可视为 $x_1$ 和 $x_2$ 的加权平均，权重分别为 $1$ 和 $\lambda$（分母归一）。
	\end{description}

	\medskip
	\textbf{\textcolor{CExplanation}{几何本质（内分与外分）}}
	$\lambda$ 的符号决定 $P$ 在线段 $P_1P_2$ 上还是在延长线上：$\lambda>0$ 时 $P$ 在 $P_1P_2$ 内部（内分点）；$\lambda<0$ 且 $\lambda\neq -1$ 时 $P$ 在外部（外分点）。当 $\lambda=1$ 时 $P$ 恰好是中点。

	\medskip
	\textbf{\textcolor{CExplanation}{代数意义（线性表示）}}
	公式表明点 $P$ 的坐标是两端点坐标的线性组合，其组合系数满足归一化条件：$\frac{1}{1+\lambda} + \frac{\lambda}{1+\lambda} = 1$。这一观点可推广到向量的线性组合。

	\medskip
	\textbf{\textcolor{CExplanation}{考点价值（坐标计算）}}
	定比分点公式是解析几何中处理共线点分点坐标的重要工具，常用于求重心坐标、线段定比分点、向量共线问题的坐标化求解。
	\begin{itemize}[leftmargin=*, itemsep=0.5em, topsep=0.4em]
		\item \textbf{考法一（直接套公式）：} 给出 $P_1,P_2$ 和 $\lambda$，直接代入公式求 $P$ 坐标。
		\item \textbf{考法二（反求参数）：} 已知分点 $P$ 坐标和端点，求 $\lambda$；或已知 $P$ 和 $\lambda$ 求某一端点。
	\end{itemize}

	\medskip
	\textbf{\textcolor{CExplanation}{顿悟点}}
	定比分点公式实质上就是“向量定比分点”的坐标表达，其核心思想是“线性加权”。记忆时可联想“$\frac{x_1+\lambda x_2}{1+\lambda}$”，分子是坐标与对应系数的乘积和，分母是系数和。

	\medskip
	\textbf{\textcolor{CExplanation}{使用场景}}
	\begin{itemize}[leftmargin=*, itemsep=0.5em, topsep=0.4em]
		\item \textbf{场景一（线段分点）：} 已知线段两端点坐标，求分点（如重心、三等分点）坐标。
		\item \textbf{场景二（向量共线）：} 在向量问题中，由三点共线条件转化为定比分点公式，实现坐标统一。
	\end{itemize}

\end{explanationbox}
